% ============================================================================
\chapter{模形式、椭圆曲线与模曲线}
\label{ch:forms-curves}
% ============================================================================
% 对应 Diamond \& Shurman Chapter 1

% --- 章导言：从一个具体函数出发，引起好奇心 ---

在正式定义之前，让我们先认识一个具体的函数。取
\[
  E_4(\tau) = 1 + 240\sum_{n=1}^{\infty} \sigma_3(n)\, q^n,
  \qquad q = e^{2\pi i\tau},
\]
其中 $\sigma_3(n) = \sum_{d \mid n} d^3$ 是因子的立方和。
前几项展开为
\[
  E_4(\tau) = 1 + 240q + 2160q^2 + 6720q^3 + \cdots
\]
这个级数收敛的区域是\textbf{上半平面} $\HH = \{\tau \in \C \mid \im(\tau) > 0\}$
（因为此时 $|q| = e^{-2\pi \im(\tau)} < 1$）。

$E_4$ 有两个令人惊叹的对称性：
\begin{enumerate}
  \item \textbf{平移不变}：$E_4(\tau + 1) = E_4(\tau)$
    ——这很"显然"，因为 $e^{2\pi i(\tau+1)} = e^{2\pi i\tau}$，
    所以 $q$ 没变。
  \item \textbf{反演对称}：$E_4(-1/\tau) = \tau^4 \cdot E_4(\tau)$
    ——这就不那么明显了！把 $\tau$ 替换成 $-1/\tau$，
    整个函数居然只差一个 $\tau^4$ 的因子。
\end{enumerate}
这些性质的证明并不平凡：反演对称性需要对求和指标做巧妙的重排
（Hecke 的技巧），而 $q$-展开的推导则用到 Poisson 求和公式。
我们将在\cref{ch:eisenstein}中给出完整证明。

这两条对称性把 $E_4$ 变成了一个非常特殊的函数，叫做\textbf{模形式}（modular form）。
它的"权"（weight）是 $4$——就是那个 $\tau^4$ 的指数。

本章的任务就是回答三个问题：
\begin{itemize}
  \item \textbf{对称性从哪来？}
    平移 $\tau \mapsto \tau+1$ 和反演 $\tau \mapsto -1/\tau$
    其实是某个群的两个生成元——它就是\textbf{模群} $\SL_2(\Z)$。
  \item \textbf{模形式空间有多大？}
    满足这些对称性的函数居然只构成一个\textbf{有限维}向量空间。
    例如权 $12$ 的模形式空间只有 $2$ 维！
  \item \textbf{模形式跟椭圆曲线有什么关系？}
    上半平面在模群作用下的商空间参数化了所有椭圆曲线。
    这个联系最终通向费马大定理。
\end{itemize}


% ============================================================================
\section{模群与上半平面}
\label{sec:modular-group-action}
% ============================================================================

上面我们观察到 $E_4$ 在平移 $\tau \mapsto \tau + 1$ 和反演 $\tau \mapsto -1/\tau$
下有特殊行为。这两个变换可以组合出更多变换，它们构成一个群。
让我们把这个群精确地描述出来。

\begin{definition}[上半平面]
\label{def:upper-half-plane}
\textbf{上半平面}\index{上半平面 (upper half-plane)}\index{H@$\mathbb{H}$}（upper half-plane）定义为
\[
  \HH = \{ \tau \in \C \mid \im(\tau) > 0 \}.
\]
\end{definition}

\begin{definition}[模群]
\label{def:modular-group}
\index{模群 (modular group)}\index{SL2Z@$\SL_2(\mathbb{Z})$}
\textbf{模群} $\SL_2(\Z)$ 是所有行列式为 $1$ 的 $2 \times 2$ 整数矩阵组成的群：
\[
  \SL_2(\Z) = \left\{
    \gamma = \begin{pmatrix} a & b \\ c & d \end{pmatrix}
    \;\middle|\; a, b, c, d \in \Z,\; ad - bc = 1
  \right\}.
\]
它通过\textbf{莫比乌斯变换}\index{莫比乌斯变换 (Möbius transformation)}（Möbius transformation）作用在 $\HH$ 上：
\[
  \gamma(\tau) = \frac{a\tau + b}{c\tau + d}.
\]
\end{definition}

\begin{proposition}
\label{prop:mobius-preserves-H}
莫比乌斯变换保持上半平面不变：若 $\tau \in \HH$，$\gamma \in \SL_2(\Z)$，
则 $\gamma(\tau) \in \HH$。具体地，
\[
  \im(\gamma(\tau)) = \frac{\im(\tau)}{|c\tau + d|^2}.
\]
\end{proposition}

\begin{proof}
设 $\gamma = \begin{pmatrix} a & b \\ c & d \end{pmatrix}$，$\tau = x + iy$。
\begin{align}
  \gamma(\tau)
  &= \frac{a\tau + b}{c\tau + d}
  = \frac{(a\tau + b)\overline{(c\tau + d)}}{|c\tau + d|^2}.
\end{align}
分子的虚部为
\begin{align}
  \im\bigl((a\tau + b)\overline{(c\tau + d)}\bigr)
  &= \im\bigl((a\tau + b)(\bar c \bar\tau + \bar d)\bigr) \\
  &= \im\bigl(ac|\tau|^2 + ad\tau + bc\bar\tau + bd\bigr) \\
  &= (ad - bc)\im(\tau) = \im(\tau).
\end{align}
最后一步用了 $ad - bc = 1$ 以及 $\im(ad\tau + bc\bar\tau) = (ad - bc)y$。
故 $\im(\gamma(\tau)) = \im(\tau)/|c\tau + d|^2 > 0$。
\end{proof}

\textbf{为什么用 $\SL_2$ 而非 $\GL_2$？}
$\GL_2(\Z)$ 包含行列式为 $\pm 1$ 的矩阵。行列式为 $-1$ 的矩阵
（如 $\begin{pmatrix} 1 & 0 \\ 0 & -1 \end{pmatrix}$）对应变换
$\tau \mapsto -\tau$，把上半平面映到\emph{下}半平面。由上面的证明，
$\im(\gamma(\tau)) = (ad - bc)\im(\tau)/|c\tau + d|^2$，
只有 $ad - bc = +1$（即 $\det \gamma = 1$）时虚部才保持正性。

\textbf{$-I$ 的作用与 $\PSL_2(\Z)$。}
$-I = \begin{pmatrix} -1 & 0 \\ 0 & -1 \end{pmatrix}$ 作用平凡：
$(-I)(\tau) = \frac{-\tau}{-1} = \tau$。
也就是说，$\SL_2(\Z)$ 中不同的矩阵可能对 $\HH$ 做同样的事
（$\gamma$ 和 $-\gamma$ 永远给出相同的变换）。
把这种"冗余"消除掉，就得到
$\PSL_2(\Z) = \SL_2(\Z)/\{\pm I\}$——
在 $\PSL_2(\Z)$ 中，不同的元素一定对应不同的变换。

\begin{theorem}[$\SL_2(\Z)$ 的生成元]
\label{thm:sl2z-generators}
$\SL_2(\Z)$ 由以下两个矩阵生成：
\[
  S = \begin{pmatrix} 0 & -1 \\ 1 & 0 \end{pmatrix}, \quad
  T = \begin{pmatrix} 1 & 1 \\ 0 & 1 \end{pmatrix}.
\]
它们对应的莫比乌斯变换为
\[
  S(\tau) = -\frac{1}{\tau}, \qquad T(\tau) = \tau + 1.
\]
注意 $S$ 在 $\SL_2(\Z)$ 中阶为 $4$（$S^2 = -I$，$S^4 = I$），
在 $\PSL_2(\Z)$ 中阶为 $2$。
\end{theorem}

\textbf{$S$ 与 $T$ 的几何直觉。}
\label{rem:S-T-geometry}
\begin{itemize}
  \item $T(\tau) = \tau + 1$：水平平移。将竖直带
    $\{-1/2 < \Re(\tau) \leq 1/2\}$ 左右"复制粘贴"铺满整个上半平面。
  \item $S(\tau) = -1/\tau$：关于单位圆的反演（加上实轴的反射）。
    关键性质：$|S(\tau)| = 1/|\tau|$，因此
    \begin{itemize}
      \item 单位圆 $|\tau| = 1$ 被映到\emph{自身}；
      \item 单位圆外部（$|\tau| > 1$）映到内部（$|S(\tau)| < 1$），反之亦然。
    \end{itemize}
    这解释了基本域 $\mathcal{F}$ 的形状：
    底边是单位圆弧（$S$ 的分界线），两侧是 $\Re(\tau) = \pm 1/2$
    （$T$ 的分界线）。
\end{itemize}

\begin{proof}
需要证明任意 $\gamma = \begin{pmatrix} a & b \\ c & d \end{pmatrix} \in \SL_2(\Z)$
可以写成 $S$ 和 $T$ 的乘积。

我们对 $|c|$ 归纳。

\textbf{基础情形 $c = 0$}：由 $ad = 1$，得 $a = d = \pm 1$。
\[
  \begin{pmatrix} 1 & b \\ 0 & 1 \end{pmatrix} = T^b, \quad
  \begin{pmatrix} -1 & -b \\ 0 & -1 \end{pmatrix} = S^2 T^{-b}.
\]
（因为 $S^2 = -I$。）

\textbf{归纳步骤}：设 $c \neq 0$。用带余除法：$a = qc + r$，$0 \leq r < |c|$。
注意
\[
  T^{-q} \gamma
  = \begin{pmatrix} 1 & -q \\ 0 & 1 \end{pmatrix}
    \begin{pmatrix} a & b \\ c & d \end{pmatrix}
  = \begin{pmatrix} a - qc & b - qd \\ c & d \end{pmatrix}
  = \begin{pmatrix} r & b-qd \\ c & d \end{pmatrix}.
\]
再左乘 $S$：
\[
  S \cdot T^{-q} \gamma
  = \begin{pmatrix} 0 & -1 \\ 1 & 0 \end{pmatrix}
    \begin{pmatrix} r & b-qd \\ c & d \end{pmatrix}
  = \begin{pmatrix} -c & -d \\ r & b-qd \end{pmatrix}.
\]
新矩阵的左下角元素为 $r$，且 $0 \leq r < |c|$。
由归纳假设，$ST^{-q}\gamma$ 可用 $S, T$ 表示，故 $\gamma$ 也可以。
\end{proof}


% ============================================================================
\section{基本域}
\label{sec:fundamental-domain}
% ============================================================================

$\SL_2(\Z)$ 包含无穷多个矩阵，每一个都把上半平面映到自身。
自然的问题是：上半平面中哪些点在 $\SL_2(\Z)$ 看来是"本质上相同的"？

这需要\textbf{群作用}\index{群作用 (group action)}（group action）的语言。群 $G$ \textbf{作用在}集合 $X$ 上，
是指一个映射 $G \times X \to X$（记 $(g, x) \mapsto g \cdot x$），满足
$e \cdot x = x$ 和 $(g_1 g_2) \cdot x = g_1 \cdot (g_2 \cdot x)$。
在我们的情况下，$G = \SL_2(\Z)$，$X = \HH$，作用就是莫比乌斯变换。

两个关键概念：
\begin{itemize}
  \item \textbf{轨道}\index{轨道 (orbit)}（orbit）：$G \cdot \tau = \{\gamma(\tau) \mid \gamma \in G\}$，
    即 $\tau$ 能被 $G$ 映到的所有点。同一轨道中的点视为"等价"的。
  \item \textbf{基本域}：从每个轨道中恰好选一个代表元，
    这些代表元组成的区域就是基本域——上半平面的一个"最小切片"。
\end{itemize}

\begin{example}[轨道是什么样的]
\label{ex:orbit-of-2i}
取 $\tau = 2i$。它的轨道 $\SL_2(\Z) \cdot (2i)$ 包含哪些点？
\begin{itemize}
  \item $T^n(2i) = 2i + n$，$n \in \Z$
    ——整条水平线 $\{\ldots, -1+2i, \; 2i, \; 1+2i, \; 2+2i, \ldots\}$。
  \item $S(2i) = -1/(2i) = i/2$
    ——跳到了更低的位置。
  \item $T(S(2i)) = i/2 + 1 = 1 + i/2$，$T^{-1}(S(2i)) = -1 + i/2$，……
  \item $S(T(2i)) = S(1+2i) = \frac{-1}{1+2i} = \frac{-(1-2i)}{5}
    = -\frac{1}{5} + \frac{2}{5}i$
    ——又跳到了别的地方。
\end{itemize}
继续组合 $S$ 和 $T$，轨道中有无穷多个点，密密麻麻地分布在上半平面中。
这些点在 $\SL_2(\Z)$ 的视角下全是"同一个点"。

基本域的作用就是：从这一大堆等价的点里，选出一个"标准代表"。
对 $2i$ 来说，它自己就在基本域中（$|\Re(2i)| = 0 < 1/2$，$|2i| = 2 > 1$），
所以 $2i$ 就是这个轨道的标准代表。
\end{example}

\begin{definition}[基本域]
\label{def:fundamental-domain}
\index{基本域 (fundamental domain)}
群 $G$ 作用在空间 $X$ 上时，一个\textbf{基本域}（fundamental domain）是
$X$ 的子集 $\mathcal{F}$，使得：
\begin{enumerate}
  \item $X = \bigcup_{\gamma \in G} \gamma(\overline{\mathcal{F}})$
    （覆盖性）；
  \item $\gamma(\mathcal{F}) \cap \mathcal{F} = \emptyset$ 对所有
    $\gamma \neq \Id$（内部不重叠）。
\end{enumerate}
\end{definition}

\begin{theorem}[$\SL_2(\Z)$ 的标准基本域]
\label{thm:fundamental-domain}
\[
  \mathcal{F} = \left\{\tau \in \HH \;\middle|\;
    |\tau| > 1,\; -\frac{1}{2} < \Re(\tau) < \frac{1}{2}
  \right\}
\]
（即开域，边界按约定选取一半）是 $\SL_2(\Z)$ 作用在 $\HH$ 上的基本域。
更精确地：$\overline{\mathcal{F}}$ 的闭包为
\[
  \overline{\mathcal{F}} = \left\{\tau \in \HH \;\middle|\;
    |\tau| \geq 1,\; -\frac{1}{2} \leq \Re(\tau) \leq \frac{1}{2}
  \right\},
\]
且 $\HH = \bigcup_{\gamma \in \SL_2(\Z)} \gamma(\overline{\mathcal{F}})$。
\end{theorem}

\begin{proof}
\textbf{覆盖性}：给定 $\tau \in \HH$，我们找 $\gamma \in \SL_2(\Z)$
使得 $\gamma(\tau) \in \overline{\mathcal{F}}$。

由\cref{prop:mobius-preserves-H}，$\im(\gamma(\tau)) = \im(\tau)/|c\tau + d|^2$。
对固定的 $\tau$，$|c\tau + d|^2 = (cx + d)^2 + c^2 y^2 \geq c^2 y^2$，
只有有限多对 $(c,d)$ 使得 $|c\tau + d| \leq 1/y$（即使虚部增大）。
因此 $\{\im(\gamma(\tau)) \mid \gamma \in \SL_2(\Z)\}$ 有最大值，
设在 $\gamma_0$ 处取到。

先用 $T^n$（$n$ 适当）使得 $\Re(\gamma_0(\tau)) \in [-1/2, 1/2]$。
设结果为 $\tau' = T^n \gamma_0(\tau)$。

若 $|\tau'| < 1$，则 $S(\tau') = -1/\tau'$ 满足
$\im(S(\tau')) = \im(\tau')/|\tau'|^2 > \im(\tau')$，
矛盾于 $\im(\tau')$ 已是最大值。故 $|\tau'| \geq 1$，
即 $\tau' \in \overline{\mathcal{F}}$。

\textbf{内部不重叠}：设 $\tau, \tau' \in \mathcal{F}$（开域内部），
$\gamma(\tau) = \tau'$，$\gamma = \begin{pmatrix} a & b \\ c & d \end{pmatrix}
\neq \pm I$。不妨设 $\im(\tau') \geq \im(\tau)$，即
$|c\tau + d|^2 \leq 1$。

由 $\tau \in \mathcal{F}$，$\im(\tau) > \sqrt{3}/2$。
若 $|c| \geq 2$，则 $|c\tau + d| \geq |c|\im(\tau) > \sqrt{3} > 1$，矛盾。
若 $|c| = 1$，则 $|c\tau + d|^2 = |\tau + d|^2 \leq 1$（当 $c = 1$），
但 $|\tau + d|^2 = (x + d)^2 + y^2$。由 $|x| < 1/2$ 和 $y > \sqrt{3}/2$，
\[
  |\tau + d|^2 \geq y^2 > 3/4.
\]
若 $d = 0$，则 $|\tau|^2 \leq 1$，但 $\tau \in \mathcal{F}$ 要求
$|\tau| > 1$，矛盾。
类似分析 $d = \pm 1$ 的情形均导致 $\tau$ 在 $\mathcal{F}$ 的边界上而非内部。
若 $c = 0$，则 $\gamma = \pm T^b$，$\tau' = \tau \pm b$，
由 $|\Re(\tau)|, |\Re(\tau')| < 1/2$ 得 $b = 0$，$\gamma = \pm I$。
\end{proof}

\textbf{基本域的形状。}
$\overline{\mathcal{F}}$ 形如一个无限高的"拱门"：底部弧线是单位圆
$|\tau| = 1$ 的一段，两侧是 $\Re(\tau) = \pm 1/2$ 的垂线。
这个形状正好对应 $S$ 和 $T$ 的分界线：$T$ 把竖直带平移，
所以左右边界由 $T$ 粘合；$S$ 关于单位圆反演，所以底部弧线由 $S$ 粘合。

\begin{example}[判断一个点是否在基本域中]
\label{ex:point-in-F}
\leavevmode
\begin{enumerate}
  \item $\tau_1 = \frac{1}{5} + 2i$：
    $|\Re(\tau_1)| = 1/5 < 1/2$，$|\tau_1|^2 = 1/25 + 4 > 1$。
    两个条件都满足，所以 $\tau_1 \in \overline{\mathcal{F}}$。

  \item $\tau_2 = 3 + i$：
    $|\Re(\tau_2)| = 3 > 1/2$，不在基本域中。
    但 $T^{-3}(\tau_2) = \tau_2 - 3 = i$，而 $i$ 在 $\overline{\mathcal{F}}$
    的边界上（$|i| = 1$）。所以 $\tau_2$ 和 $i$ 属于同一个轨道。

  \item $\tau_3 = i/3$：
    $|\tau_3| = 1/3 < 1$，不在基本域中。
    用 $S$：$S(\tau_3) = -1/(i/3) = -3/i = 3i$。
    而 $3i$ 满足 $|\Re(3i)| = 0 < 1/2$，$|3i| = 3 > 1$，
    所以 $3i \in \mathcal{F}$。
\end{enumerate}
\end{example}

\begin{example}[把一个点"折回"基本域]
\label{ex:reduce-to-F}
取 $\tau = \frac{7}{3} + \frac{i}{2}$。怎样找到它在基本域中的代表元？

\textbf{第一步}：用 $T$ 调整实部。$\Re(\tau) = 7/3 \approx 2.33$，
取 $T^{-2}$：$\tau' = \tau - 2 = \frac{1}{3} + \frac{i}{2}$。
现在 $|\Re(\tau')| = 1/3 < 1/2$，实部条件满足了。

\textbf{第二步}：检查模长。$|\tau'|^2 = 1/9 + 1/4 = 13/36 < 1$，
不满足 $|\tau| \geq 1$。用 $S$：
\[
  S(\tau') = \frac{-1}{\frac{1}{3} + \frac{i}{2}}
  = \frac{-1}{\frac{2 + 3i}{6}} = \frac{-6}{2 + 3i}
  = \frac{-6(2 - 3i)}{4 + 9} = \frac{-12 + 18i}{13}.
\]
所以 $S(\tau') = -\frac{12}{13} + \frac{18}{13}i$。

\textbf{第三步}：再用 $T$ 调整实部。$\Re(S(\tau')) = -12/13 \approx -0.92$，
取 $T^1$：$\tau'' = S(\tau') + 1 = \frac{1}{13} + \frac{18}{13}i$。
现在 $|\Re(\tau'')| = 1/13 < 1/2$，
$|\tau''|^2 = 1/169 + 324/169 = 325/169 = 25/13 > 1$。两个条件都满足！

结论：$\tau'' = \frac{1}{13} + \frac{18}{13}i \in \mathcal{F}$，
由 $\tau'' = TS T^{-2}(\tau)$ 得到。
\end{example}

\textbf{椭圆不动点与稳定子。}
\label{rem:elliptic-fixed-points}
大多数 $\tau \in \HH$ 只有"平凡"的对称性——只有 $\pm I$
把它映回自身。但基本域边界上有两个特殊点，它们有额外的对称性。

给定 $\tau$，所有把 $\tau$ 固定不动的矩阵构成一个子群，
叫做 $\tau$ 的\textbf{稳定子}\index{稳定子 (stabilizer)}（stabilizer）：
$(\SL_2(\Z))_\tau = \{\gamma \in \SL_2(\Z) \mid \gamma(\tau) = \tau\}$。

\begin{enumerate}
  \item $\tau = i$：$S(i) = -1/i = i$，所以 $S$ 固定 $i$。
    稳定子为 $\{I, S, -I, -S\}$，同构于 $\Z/4\Z$
    （在 $\PSL_2(\Z)$ 中为 $\Z/2\Z$）。
  \item $\tau = \rho = e^{2\pi i/3} = \frac{-1+i\sqrt{3}}{2}$：
    直接计算 $(ST)(\rho) = S(\rho + 1) = \rho$。
    稳定子为 $\langle ST \rangle \cong \Z/6\Z$
    （在 $\PSL_2(\Z)$ 中为 $\Z/3\Z$）。
\end{enumerate}
这两个\textbf{椭圆不动点}导致基本域边界上的粘合出现特殊行为，
在\cref{ch:dimension}的维数公式中会产生修正项。


% ============================================================================
\section{尖点}
\label{sec:cusps}
% ============================================================================

基本域 $\mathcal{F}$ 向上延伸到无穷远——它的"顶部"越来越窄，
像一个无限高的烟囱。这个"烟囱口"在数学上需要一个名字，
它就是\textbf{尖点}（cusp）。

\textbf{为什么需要尖点？}
回忆模形式 $E_4(\tau)$ 的 $q$-展开 $\sum a_n q^n$，其中 $q = e^{2\pi i\tau}$。
当 $\im(\tau) \to +\infty$ 时，$q \to 0$。
所以"$\tau$ 趋向无穷远"对应"$q$ 趋向 $0$"，而我们要求模形式在 $q = 0$
处有良好的行为（全纯或消失）。
这个"无穷远点"就是最基本的尖点 $\infty$。

\textbf{$\mathbb{P}^1(\Q)$ 是什么？}
除了 $\infty$ 之外，还有其他"方向"可以趋向实轴上的有理点。
$\mathbb{P}^1(\Q)$ 就是 $\Q \cup \{\infty\}$——
所有有理数加上一个无穷远点，可以理解为"有理数的圆周"。

\begin{definition}[扩展上半平面与尖点]
\label{def:cusps}
\index{扩展上半平面}\index{尖点 (cusp)}\index{P1Q@$\mathbb{P}^1(\mathbb{Q})$}
\textbf{扩展上半平面}定义为
\[
  \HH^* = \HH \cup \Q \cup \{\infty\} = \HH \cup \mathbb{P}^1(\Q).
\]
$\SL_2(\Z)$ 的莫比乌斯变换自然延拓到 $\mathbb{P}^1(\Q)$ 上：
对 $\gamma = \begin{pmatrix} a & b \\ c & d \end{pmatrix}$，
\[
  \gamma(\infty) = \frac{a}{c}, \qquad
  \gamma\!\left(-\frac{d}{c}\right) = \infty.
\]
（就是把 $\frac{a\tau + b}{c\tau + d}$ 中分母为零和分子为零的情况补上。）

$\mathbb{P}^1(\Q)$ 中的点称为\textbf{尖点}（cusps）。
\end{definition}

\begin{example}[具体的尖点变换]
\label{ex:cusp-action}
\leavevmode
\begin{itemize}
  \item $T(\infty) = \infty + 1 = \infty$
    ——平移不影响无穷远点。
  \item $S(\infty) = -1/\infty = 0$
    ——反演把 $\infty$ 映到 $0$。
  \item $S(0) = -1/0 = \infty$
    ——反演把 $0$ 映回 $\infty$。
  \item 取 $\gamma = \begin{pmatrix} 3 & 1 \\ 2 & 1 \end{pmatrix}
    \in \SL_2(\Z)$（$3 \cdot 1 - 1 \cdot 2 = 1$），
    则 $\gamma(\infty) = 3/2$。
    所以 $3/2$ 和 $\infty$ 在同一轨道中。
\end{itemize}
\end{example}

\begin{proposition}
\label{prop:cusps-transitive}
$\SL_2(\Z)$ 在 $\mathbb{P}^1(\Q)$ 上的作用是传递的，即所有尖点都等价。
换句话说：$\SL_2(\Z) \backslash \mathbb{P}^1(\Q)$ 只有一个点。
\end{proposition}

\begin{proof}
任取 $a/c \in \Q$（$\gcd(a,c) = 1$），由 Bézout 恒等式，存在 $b, d \in \Z$
使得 $ad - bc = 1$。则
$\gamma = \begin{pmatrix} a & b \\ c & d \end{pmatrix} \in \SL_2(\Z)$
满足 $\gamma(\infty) = a/c$。所以每个有理数都和 $\infty$ 在同一轨道中。
\end{proof}

\textbf{同余子群的尖点。}
对 $\SL_2(\Z)$ 来说只有一个尖点等价类，但对更小的同余子群
（如 $\GammaO(N)$），尖点会分裂成多个等价类。
例如 $\GammaO(2)$ 有 $2$ 个不等价的尖点：$\infty$ 和 $0$
（因为 $S \notin \GammaO(2)$，无法把 $\infty$ 映到 $0$）。
这将在\cref{sec:congruence-subgroups}中进一步讨论。


% ============================================================================
\section{模形式的定义}
\label{sec:modular-forms-def}
% ============================================================================

在章开头，我们已经看到 $E_4(\tau)$ 满足 $E_4(\tau+1) = E_4(\tau)$ 和
$E_4(-1/\tau) = \tau^4 E_4(\tau)$。现在我们把这种对称性抽象为一般定义。

\begin{definition}[权 $k$ 的弱模形式]
\label{def:weakly-modular}
\index{弱模形式 (weakly modular form)}
设 $k \in \Z$。全纯函数 $f\colon \HH \to \C$ 称为权 $k$ 的
\textbf{弱模形式}（weakly modular form），若对所有
$\gamma = \begin{pmatrix} a & b \\ c & d \end{pmatrix} \in \SL_2(\Z)$，
\[
  f(\gamma(\tau)) = (c\tau + d)^k f(\tau).
\]
由\cref{thm:sl2z-generators}，这等价于同时满足：
\begin{enumerate}
  \item $f(\tau + 1) = f(\tau)$（由 $T$ 给出）；
  \item $f(-1/\tau) = \tau^k f(\tau)$（由 $S$ 给出）。
\end{enumerate}
\end{definition}

条件 (1) 意味着 $f$ 有以 $q = e^{2\pi i\tau}$ 为变量的 Fourier 展开，
称为 \textbf{$q$-展开}\index{q-展开 ($q$-expansion)}（$q$-expansion）：
\[
  f(\tau) = \sum_{n \in \Z} a_n q^n, \quad q = e^{2\pi i\tau}.
\]

\begin{definition}[模形式与尖点形式]
\label{def:modular-form}
\index{模形式 (modular form)}\index{尖点形式 (cusp form)}\index{Mk@$\mathcal{M}_k$}\index{Sk@$\mathcal{S}_k$}
权 $k$ 的\textbf{模形式}（modular form）是权 $k$ 的弱模形式 $f$，
且 $f$ 在尖点 $\infty$ 处\textbf{全纯}，即 $q$-展开中 $a_n = 0$ 对所有 $n < 0$：
\[
  f(\tau) = \sum_{n=0}^{\infty} a_n q^n.
\]
权 $k$ 的模形式空间记为 $\Mk_k(\SL_2(\Z))$。

若进一步 $a_0 = 0$（即 $f$ 在尖点处消失），则 $f$ 称为
\textbf{尖点形式}（cusp form）。尖点形式空间记为 $\Sk_k(\SL_2(\Z))$。
\end{definition}

\begin{example}[Eisenstein 级数]
\label{ex:eisenstein-preview}
对偶整数 $k \geq 4$，\textbf{Eisenstein 级数}
\[
  G_k(\tau) = \sum_{\substack{(c,d) \in \Z^2 \\ (c,d) \neq (0,0)}}
  \frac{1}{(c\tau + d)^k}
\]
是权 $k$ 的模形式。我们将在\cref{ch:eisenstein}中详细讨论。
规范化版本 $E_k = G_k / (2\zeta(k))$ 的 $q$-展开为
\[
  E_k(\tau) = 1 - \frac{2k}{B_k} \sum_{n=1}^{\infty} \sigma_{k-1}(n) q^n,
\]
其中 $B_k$ 是 Bernoulli 数，$\sigma_{k-1}(n) = \sum_{d \mid n} d^{k-1}$。
\end{example}

\begin{example}[判别形式 $\Delta$]
\label{ex:discriminant-form}
\textbf{判别形式}
\[
  \Delta(\tau) = \frac{1}{1728}\bigl(E_4(\tau)^3 - E_6(\tau)^2\bigr)
  = q \prod_{n=1}^{\infty}(1 - q^n)^{24}
  = \sum_{n=1}^{\infty} \tau(n) q^n
\]
是权 $12$ 的尖点形式，其中 $\tau(n)$ 是 \textbf{Ramanujan $\tau$-函数}。
$\Delta$ 在 $\HH$ 上无零点。
\end{example}

\begin{example}[$j$-不变量]
\label{ex:j-invariant-function}
\textbf{$j$-不变量}
\[
  j(\tau) = \frac{E_4(\tau)^3}{\Delta(\tau)} = \frac{1}{q} + 744 + 196884q + \cdots
\]
是 $\SL_2(\Z)$ 的一个\textbf{模函数}（权 $0$ 的亚纯模形式），
它建立了 $\SL_2(\Z) \backslash \HH$ 到 $\C$ 的双射。
\end{example}


% ============================================================================
\section{同余子群}
\label{sec:congruence-subgroups}
% ============================================================================

\begin{definition}[同余子群]
\label{def:congruence-subgroups}
\index{同余子群 (congruence subgroup)}\index{Gamma0@$\Gamma_0(N)$}\index{Gamma1@$\Gamma_1(N)$}\index{GammaN@$\Gamma(N)$}
设 $N$ 为正整数。定义以下 $\SL_2(\Z)$ 的子群：
\begin{align}
  \Gamma(N) &= \left\{
    \begin{pmatrix} a & b \\ c & d \end{pmatrix} \in \SL_2(\Z)
    \;\middle|\;
    \begin{pmatrix} a & b \\ c & d \end{pmatrix}
    \equiv \begin{pmatrix} 1 & 0 \\ 0 & 1 \end{pmatrix} \pmod{N}
  \right\}, \\[6pt]
  \GammaI(N) &= \left\{
    \begin{pmatrix} a & b \\ c & d \end{pmatrix} \in \SL_2(\Z)
    \;\middle|\;
    a \equiv d \equiv 1,\; c \equiv 0 \pmod{N}
  \right\}, \\[6pt]
  \GammaO(N) &= \left\{
    \begin{pmatrix} a & b \\ c & d \end{pmatrix} \in \SL_2(\Z)
    \;\middle|\;
    c \equiv 0 \pmod{N}
  \right\}.
\end{align}
包含关系为 $\Gamma(N) \trianglelefteq \GammaI(N) \trianglelefteq \GammaO(N)
\leq \SL_2(\Z)$。

更一般地，$\SL_2(\Z)$ 的子群 $\Gamma$ 若包含某个 $\Gamma(N)$，
则称为\textbf{同余子群}（congruence subgroup），满足此条件的最小 $N$ 称为
$\Gamma$ 的\textbf{级}（level）。
\end{definition}

\begin{proposition}[指标公式]
\label{prop:index-formulas}
\leavevmode
\begin{enumerate}
  \item $[\SL_2(\Z) : \Gamma(N)] = N^3 \prod_{p \mid N}(1 - 1/p^2)$
    （$N \geq 1$）。
  \item $[\SL_2(\Z) : \GammaI(N)] = N^2 \prod_{p \mid N}(1 - 1/p^2)$
    （$N \geq 1$）。
  \item $[\SL_2(\Z) : \GammaO(N)] = N \prod_{p \mid N}(1 + 1/p)$
    （$N \geq 1$）。
\end{enumerate}
\end{proposition}

\begin{proof}
考虑约化同态 $\SL_2(\Z) \to \SL_2(\Z/N\Z)$。

\textbf{(1)}：$\Gamma(N) = \ker(\SL_2(\Z) \to \SL_2(\Z/N\Z))$。
此约化同态是满射（这需要证明，我们在习题中给出提示）。
故 $[\SL_2(\Z) : \Gamma(N)] = \#\SL_2(\Z/N\Z)$。
由中国剩余定理，$\#\SL_2(\Z/N\Z) = \prod_{p^e \| N} \#\SL_2(\Z/p^e\Z)$。
而 $\#\SL_2(\Z/p^e\Z) = p^{3e} \prod_{p \mid p^e}(1 - 1/p^2)
= p^{3(e-1)} \cdot p^3(1 - 1/p^2) = p^{3e}(1 - 1/p^2)$。

\textbf{(2)}：$\GammaI(N)/\Gamma(N) \cong \Z/N\Z$
（通过映射 $\gamma \mapsto b \bmod N$），故
$[\SL_2(\Z) : \GammaI(N)] = [\SL_2(\Z) : \Gamma(N)]/N$。

\textbf{(3)}：$\GammaO(N)/\GammaI(N) \cong (\Z/N\Z)^\times$
（通过映射 $\gamma \mapsto d \bmod N$），故
$[\SL_2(\Z) : \GammaO(N)] = [\SL_2(\Z) : \GammaI(N)]/\varphi(N)$。
化简得 $N \prod_{p \mid N}(1 + 1/p)$。
\end{proof}

\begin{definition}[同余子群上的模形式]
\label{def:modular-forms-congruence}
设 $\Gamma$ 为同余子群。权 $k$ 的\textbf{模形式} $f \in \Mk_k(\Gamma)$
是全纯函数 $f\colon \HH \to \C$，满足：
\begin{enumerate}
  \item 对所有 $\gamma = \begin{pmatrix} a & b \\ c & d \end{pmatrix}
    \in \Gamma$，
    \[
      f(\gamma(\tau)) = (c\tau + d)^k f(\tau);
    \]
  \item $f$ 在 $\Gamma$ 的每个尖点处全纯。
\end{enumerate}
若 $f$ 在所有尖点处消失，则 $f$ 是\textbf{尖点形式}，$f \in \Sk_k(\Gamma)$。
\end{definition}


% ============================================================================
\section{复环面}
\label{sec:complex-tori}
% ============================================================================

\begin{definition}[格与复环面]
\label{def:lattice-torus}
\index{格 (lattice)}\index{复环面 (complex torus)}
$\C$ 中的一个\textbf{格}（lattice）是
$\Lambda = \Z\omega_1 + \Z\omega_2$（$\omega_1/\omega_2 \notin \R$）。
对应的\textbf{复环面}为商群 $\C/\Lambda$，它是一个亏格 $1$ 的紧黎曼面。
\end{definition}

\begin{proposition}[复环面的同构分类]
\label{prop:torus-isomorphism}
$\C/\Lambda \cong \C/\Lambda'$（作为复 Lie 群）当且仅当
$\Lambda' = \alpha\Lambda$ 对某 $\alpha \in \C^*$。
\end{proposition}

\begin{proof}
设 $\varphi\colon \C/\Lambda \to \C/\Lambda'$ 为全纯群同构。
提升到万有覆盖 $\C$，得到全纯函数 $\tilde\varphi\colon \C \to \C$，
满足 $\tilde\varphi(z + \omega) - \tilde\varphi(z) \in \Lambda'$ 对所有
$\omega \in \Lambda$。

由于 $z \mapsto \tilde\varphi(z + \omega) - \tilde\varphi(z)$ 连续且值在离散集
$\Lambda'$ 中，它必为常数。对 $\omega$ 求导得
$\tilde\varphi'(z + \omega) = \tilde\varphi'(z)$，即 $\tilde\varphi'$ 以 $\Lambda$
为周期。由\cref{thm:elliptic-basic}(1)（全纯椭圆函数为常数），
$\tilde\varphi'(z) = \alpha$（常数）。故 $\tilde\varphi(z) = \alpha z + \beta$。
由 $\tilde\varphi(\Lambda) \subset \Lambda'$ 和满射性，$\alpha\Lambda = \Lambda'$。
\end{proof}

\begin{corollary}
\label{cor:torus-tau}
每个复环面同构于 $\C/\Lambda_\tau$（$\Lambda_\tau = \Z\tau + \Z$，
$\tau \in \HH$）。两个环面 $\C/\Lambda_\tau$ 和 $\C/\Lambda_{\tau'}$
同构当且仅当 $\tau' = \gamma(\tau)$ 对某 $\gamma \in \SL_2(\Z)$。

因此，复环面的同构类由 $\SL_2(\Z) \backslash \HH$ 的点参数化。
\end{corollary}

\begin{proof}
设 $\Lambda = \Z\omega_1 + \Z\omega_2$，$\im(\omega_1/\omega_2) > 0$。
取 $\alpha = 1/\omega_2$，则 $\alpha\Lambda = \Z(\omega_1/\omega_2) + \Z
= \Lambda_\tau$（$\tau = \omega_1/\omega_2 \in \HH$）。

$\C/\Lambda_\tau \cong \C/\Lambda_{\tau'}$ 当且仅当
$\Lambda_{\tau'} = \alpha \Lambda_\tau$ 对某 $\alpha \in \C^*$。
设 $\alpha \tau = a\tau' + b$，$\alpha = c\tau' + d$（$a,b,c,d \in \Z$），
则 $\tau = (a\tau' + b)/(c\tau' + d)$。由 $\im(\tau) > 0$、$\im(\tau') > 0$，
要求 $ad - bc = 1$。
\end{proof}


% ============================================================================
\section{复环面与椭圆曲线}
\label{sec:tori-and-ec}
% ============================================================================

\begin{theorem}[均匀化定理]
\label{thm:uniformization}
在 $\C$ 上，复环面 $\C/\Lambda$ 与椭圆曲线之间有经典的一一对应：

\textbf{从环面到曲线}：Weierstrass $\wp$ 函数给出同构
\[
  \C/\Lambda \xrightarrow{\;\sim\;} E_\Lambda(\C), \quad
  z \mapsto [\wp(z;\Lambda) : \wp'(z;\Lambda) : 1],
\]
其中 $E_\Lambda\colon Y^2 = 4X^3 - g_2(\Lambda)X - g_3(\Lambda)$，
$g_2 = 60G_4(\Lambda)$，$g_3 = 140G_6(\Lambda)$。

\textbf{从曲线到环面}：给定 $\C$ 上的椭圆曲线 $E$，存在唯一的格 $\Lambda$
（在相似变换下）使得 $E(\C) \cong \C/\Lambda$。
\end{theorem}

\begin{proof}
第一个方向已在\cref{thm:weierstrass-properties}中证明。

第二个方向：设 $E\colon y^2 = 4x^3 - Ax - B$（$A^3 - 27B^2 \neq 0$）。
需要找格 $\Lambda$ 使得 $g_2(\Lambda) = A$，$g_3(\Lambda) = B$。
这等价于要求 $j(\Lambda) = 1728 A^3/(A^3 - 27B^2)$。

$j\colon \SL_2(\Z) \backslash \HH \to \C$ 是双射（\cref{ex:j-invariant-function}），
故存在唯一的 $\tau \in \mathcal{F}$ 使得 $j(\tau)$ 等于所需值。
然后对 $\Lambda_\tau$ 做适当缩放使 $g_2, g_3$ 匹配。
\end{proof}

\begin{definition}[同源]
\label{def:isogeny}
\index{同源 (isogeny)}
设 $E_1, E_2$ 为椭圆曲线。一个\textbf{同源}（isogeny）
$\varphi\colon E_1 \to E_2$ 是满足 $\varphi(O) = O$ 的代数群态射。

在复环面语言中，$\C/\Lambda_1 \to \C/\Lambda_2$ 的同源对应于
$\alpha \in \C^*$ 使得 $\alpha\Lambda_1 \subset \Lambda_2$，
映射为 $z \mapsto \alpha z$。同源的\textbf{度}为 $[\Lambda_2 : \alpha\Lambda_1]$。
\end{definition}


% ============================================================================
\section{模曲线与模空间}
\label{sec:modular-curves-moduli}
% ============================================================================

\begin{definition}[模曲线]
\label{def:modular-curves}
对同余子群 $\Gamma$，定义\textbf{（开）模曲线}
\[
  Y(\Gamma) = \Gamma \backslash \HH.
\]
特别地，$Y_0(N) = \GammaO(N) \backslash \HH$，
$Y_1(N) = \GammaI(N) \backslash \HH$，
$Y(N) = \Gamma(N) \backslash \HH$。

通过添加尖点得到\textbf{（紧化）模曲线}
\[
  X(\Gamma) = \Gamma \backslash \HH^*.
\]
$X(\Gamma)$ 是紧黎曼面（将在\cref{ch:riemann-surfaces}中详细讨论）。
\end{definition}

\begin{theorem}[模诠释]
\label{thm:moduli-interpretation}
模曲线参数化带附加结构的椭圆曲线（在 $\C$ 上）：
\begin{enumerate}
  \item $Y(1) = \SL_2(\Z) \backslash \HH$ 参数化复环面的同构类
    （\cref{cor:torus-tau}）。通过 $j$-不变量，$Y(1) \cong \C$。
  \item $Y_0(N)$ 参数化 $\{(E, C)\}$ 的同构类，其中 $E$ 是椭圆曲线，
    $C \subset E$ 是 $N$ 阶循环子群（等价于核为 $\Z/N\Z$ 的同源
    $E \to E/C$）。
  \item $Y_1(N)$ 参数化 $\{(E, P)\}$ 的同构类，其中 $P \in E$ 是
    $N$ 阶点。
  \item $Y(N)$ 参数化 $\{(E, (P, Q))\}$ 的同构类，其中 $(P, Q)$ 是
    $E[N]$ 的一组有序基。
\end{enumerate}
\end{theorem}

\begin{proof}[解释]
\textbf{(2)}：设 $\tau \in \HH$，对应环面 $E_\tau = \C/\Lambda_\tau$
（$\Lambda_\tau = \Z\tau + \Z$）。取 $C = \langle 1/N \rangle \subset E_\tau$，
即 $C = \frac{1}{N}\Z/\Z$ 是 $N$ 阶循环子群。

何时 $(E_\tau, C_\tau) \cong (E_{\tau'}, C_{\tau'})$？
需要 $\alpha \in \C^*$ 使得 $\alpha\Lambda_\tau = \Lambda_{\tau'}$
且 $\alpha C_\tau = C_{\tau'}$。写出条件：
$\alpha\tau = a\tau' + b$，$\alpha = c\tau' + d$（$ad - bc = 1$），
且 $\alpha \cdot \frac{1}{N} \equiv \frac{j}{N} \pmod{\Lambda_{\tau'}}$
对某 $j$ 与 $N$ 互素。

这等价于 $\frac{c\tau' + d}{N} \in \frac{1}{N}\Z + \Lambda_{\tau'}$，
即 $c \equiv 0 \pmod{N}$。
故等价关系正是 $\GammaO(N)$ 的作用。

\textbf{(3) 和 (4)}：类似分析。
\end{proof}


% ============================================================================
\section{习题}
\label{sec:ch1-exercises}
% ============================================================================

\begin{exercise}
\label{ex:ch1-sl2z-1}
验证 $S^2 = -I$，$(ST)^3 = -I$。
由此说明 $\PSL_2(\Z) = \SL_2(\Z)/\{\pm I\}$ 的表现为
$\langle \bar{S}, \bar{T} \mid \bar{S}^2 = (\bar{S}\bar{T})^3 = 1 \rangle$。
\end{exercise}

\begin{exercise}
\label{ex:ch1-fund-domain}
画出 $\GammaO(2)$ 的基本域。
（提示：$[\SL_2(\Z) : \GammaO(2)] = 3$，基本域是 $\SL_2(\Z)$ 基本域的
$3$ 个拷贝的并。）
\end{exercise}

\begin{exercise}
\label{ex:ch1-cusps-gamma0}
证明 $\GammaO(N)$ 在 $\mathbb{P}^1(\Q)$ 上的作用不再传递（当 $N > 1$）。
确定 $\GammaO(N)$ 的尖点等价类个数（用 $N$ 的因子表达）。
\end{exercise}

\begin{exercise}
\label{ex:ch1-weight-k}
证明：若 $k$ 为奇数，则 $\Mk_k(\SL_2(\Z)) = 0$。
（提示：用 $-I \in \SL_2(\Z)$ 的变换律。）
\end{exercise}

\begin{exercise}
\label{ex:ch1-surjection}
证明约化同态 $\SL_2(\Z) \to \SL_2(\Z/N\Z)$ 是满射。
（提示：先证 $N = p$ 素数的情形，利用 Hensel 引理或中国剩余定理。）
\end{exercise}

\begin{exercise}
\label{ex:ch1-isogeny-degree}
设 $\Lambda = \Z\tau + \Z$。证明：$\C/\Lambda$ 到自身的度为 $n$ 的同源
恰好对应于满足 $\alpha\Lambda \subset \Lambda$、$[\Lambda : \alpha\Lambda] = n$
的 $\alpha \in \C$。当 $\Lambda$ 没有复乘（CM）时，说明这样的 $\alpha$ 必为整数。
\end{exercise}
